\documentclass[11pt,a4paper]{article}
\usepackage[utf8]{inputenc}
\usepackage[T1]{fontenc}
\usepackage{lmodern}
\usepackage[spanish]{babel}
\usepackage[margin=2.54cm]{geometry}
\usepackage{enumitem}
\usepackage{xcolor}
\usepackage{titlesec}
\usepackage{hyperref}
\usepackage[table]{xcolor}

% Configuración visual de títulos
\titleformat{\section}{\Large\bfseries\color{blue!60!black}}{\thesection.}{1em}{}
\titleformat{\subsection}{\large\bfseries\color{black!80}}{\thesubsection.}{1em}{}
\titleformat{\subsubsection}{\normalsize\bfseries}{\thesubsubsection.}{1em}{}

% Comandos personalizados para Objetivos y Resultados solicitados por el profesor
\newcommand{\objetivo}[1]{\noindent\textbf{\textcolor{blue!60!black}{Objetivo:}} #1\par\vspace{0.1cm}}
\newcommand{\resultado}[1]{\noindent\textbf{\textcolor{green!50!black}{Resultado esperado:}} #1\par\vspace{0.4cm}}

\begin{document}

% --------------------------------------------------------
% PORTADA (Corrigiendo: "No hay título del proyecto")
% --------------------------------------------------------
\begin{titlepage}
    \centering
    \vspace*{3cm}
    {\Huge \textbf{Plan de Investigación: \\ Industria Pesquera y de Procesamiento de Productos del Mar} \par}
    \vspace{1.5cm}
    {\Large \textbf{Asignatura:} Investigación por Muestreo \par}
    \vspace{0.5cm}
    {\Large \textbf{Autores:} Sandro Guevara \par}
    \vspace{2.5cm}
  
    \vfill
\end{titlepage}

\tableofcontents
\newpage

% --------------------------------------------------------
% CAPÍTULO 1
% --------------------------------------------------------
\section{Descripción de la Industria Pesquera}
\objetivo{Investigar y categorizar los fundamentos operativos, logísticos y comerciales que rigen a las industrias de captura y procesamiento de productos del mar a nivel general.}
\resultado{Un perfil detallado del sector que identifique su propósito, oferta comercial, entorno competitivo y los actores que intervienen en su cadena de valor.}

\subsection{Características Generales de la Industria}
\objetivo{Investigar y delimitar el marco institucional, productivo y estratégico que define la identidad operativa de las plantas industriales pesqueras.}
\resultado{Un diagnóstico situacional estructurado acerca de la naturaleza de la industria, sus productos principales y el comportamiento de su mercado.}

\subsubsection{Naturaleza, Razón y Propósito}
La industria pesquera tiene como propósito fundamental la extracción de recursos hidrobiológicos (pesca extractiva) y su posterior transformación termo-mecánica o química para el consumo humano e industrial. Su razón de ser radica en garantizar el suministro de proteína marina de alta calidad a mercados globales, así como proveer materias primas esenciales (harina y aceite de pescado) para sustentar la cadena alimenticia del sector agropecuario y acuícola a nivel mundial.

\subsubsection{Productos y Líneas de Producción}
La investigación levantará información sobre las principales líneas de salida de estas plantas industriales, típicamente ramificadas en:
\begin{itemize}[noitemsep]
    \item \textbf{Línea de Congelados (IQF):} Pescado entero, lomos o filetes sometidos a túneles de congelación rápida individual para exportación marítima.
    \item \textbf{Línea de Conservas:} Productos cárnicos marinos cocidos y sellados herméticamente (enlatados y \textit{pouches}) con tratamientos de esterilización que garantizan la inocuidad sin necesidad de cadena de frío.
    \item \textbf{Línea de Reducción (Consumo indirecto):} Subproductos derivados del fileteo y especies pelágicas menores transformadas en harina y aceite industrial.
\end{itemize}

\subsubsection{Clientes y Canales de Distribución}
Se investigará la segmentación de la demanda de esta industria:
\begin{itemize}[noitemsep]
    \item \textbf{Mercado B2B (Industrial y Mayorista):} Fábricas de alimento balanceado, marcas internacionales que subcontratan procesos (maquila), y grandes mayoristas que adquieren el producto por contenedores.
    \item \textbf{Mercado B2C (Consumo Final a través de \textit{Retail}):} Cadenas de supermercados que exigen altos estándares de presentación, gramaje exacto y etiquetado nutricional para distribución directa al consumidor.
\end{itemize}

\subsubsection{Estructura de Competencia y Proveedores}
Se indagará en la vulnerabilidad de la cadena de suministro y la fuerte regulación competitiva del sector:
\begin{itemize}[noitemsep]
    \item \textbf{Proveedores Críticos:} Astilleros navales (mantenimiento de flota), industrias petroquímicas (combustible y amoníaco), fábricas de envases de hojalata y proveedores de redes de cerco.
    \item \textbf{Competencia y Entorno Legal:} La competencia no solo está dictada por precios, sino por el acceso a la materia prima limitado por \textbf{cuotas de captura} y \textbf{vedas biológicas} dictaminadas por las autoridades gubernamentales. Se investigará el cumplimiento de normativas de salubridad y trazabilidad internacional.
\end{itemize}

% --------------------------------------------------------
% CAPÍTULO 2
% --------------------------------------------------------
\section{Procesos Sistémicos de la Industria }
\objetivo{Descomponer la cadena de valor de la industria en sus fases de transformación para entender cómo fluyen la materia prima y la información a través del sistema.}
\resultado{Un mapa técnico descriptivo que identifique las entradas, las fases de transformación, las salidas y los mecanismos de retroalimentación de la planta industrial.}

\subsection{El Proceso de Transformación}
\objetivo{Mapear analíticamente la secuencia física que atraviesan los recursos marinos desde su descarga en muelle hasta su consolidación en bodega térmica.}
\resultado{Un esquema teórico del flujo operativo estándar, identificando los puntos donde la materia prima sufre cambios físicos, químicos o térmicos.}

\subsubsection{Entradas y Recursos }
Se investigará la matriz de insumos sin los cuales la industria detiene sus operaciones:
\begin{itemize}[noitemsep]
    \item \textbf{Materia Prima Biológica:} Captura marina fresca o congelada en salmuera. Su calidad de entrada condiciona todo el sistema.
    \item \textbf{Insumos Físicos e Industriales:} Hielo en escamas, volúmenes masivos de agua purificada, y materiales de empaque con grado alimenticio.
    \item \textbf{Recursos Energéticos:} Energía eléctrica trifásica indispensable para sostener las cámaras frigoríficas y vapor de calderas usado en los procesos térmicos.
\end{itemize}

\subsubsection{Fases de Producción }
El plan investigará cómo ocurren las transformaciones clave dentro de la nave industrial:
\begin{enumerate}[noitemsep]
    \item \textbf{Recepción y Clasificación:} Descarga controlada, lavado, separación por tallas y evaluación organoléptica inicial antes del ingreso a planta.
    \item \textbf{Acondicionamiento (Eviscerado y Corte):} Retiro de elementos no comerciales (vísceras, espinas, escamas) mediante procesos mecánicos o manuales en ambientes climatizados para evitar proliferación bacteriana.
    \item \textbf{Transformación Térmica:} Aplicación de curvas de temperatura. Dependiendo de la línea, la carne ingresa a cocedores de vapor continuo o a túneles de congelación extrema.
    \item \textbf{Empaque y Sellado:} Inserción del producto procesado en sus contenedores finales, garantizando el peso drenado y el vacío hermético para su protección oxidativa.
\end{enumerate}

\subsubsection{Salidas y Retroalimentación }
Se estudiará cómo la industria mide su éxito y utiliza la información para corregir errores:
\begin{itemize}[noitemsep]
    \item \textbf{Salidas Físicas:} Lotes de productos aptos para consumo, subproductos aprovechables, y efluentes industriales que requieren tratamiento de aguas.
    \item \textbf{Bucle de Retroalimentación Interno:} Mecanismos de comunicación en planta. Por ejemplo, si control de calidad detecta mermas excesivas en el área de corte, se retroalimenta inmediatamente a los operarios para ajustar la técnica.
    \item \textbf{Bucle de Retroalimentación Externo:} Evaluación post-venta. Reclamos de clientes por rechazos en puerto (ej. pérdida de cadena de frío) obligan a la industria a auditar retrospectivamente a sus empresas de transporte naviero.
\end{itemize}

\subsection{Fases de Ejecución del Plan}
\objetivo{Estructurar una hoja de ruta metodológica y no invasiva para la recopilación exhaustiva de datos operativos, comerciales y normativos del sector pesquero.}
\resultado{Un cronograma de investigación puramente analítico que culmine en la identificación y parametrización de las variables críticas del sistema industrial.}

Para comprender a la industria en su totalidad sin intervenir en sus operaciones, el equipo ejecutará las siguientes fases de investigación pura:

\begin{itemize}
    \item \textbf{Fase 1: Revisión del Ecosistema Normativo y Estructural.}
    \begin{itemize}[noitemsep, label=-]
        \item Identificación de regulaciones gubernamentales restrictivas (leyes de pesca, períodos de veda, cuotas de captura por especie).
        \item Auditoría teórica de estándares internacionales de calidad y trazabilidad exigidos para exportación (HACCP, ISO 22000, normativas de la FAO).
        \item Análisis de organigramas tipo de la industria para mapear la jerarquía departamental (flota, planta, mantenimiento, aseguramiento de calidad).
    \end{itemize}

    \item \textbf{Fase 2: Levantamiento Analítico del Proceso de Transformación.}
    \begin{itemize}[noitemsep, label=-]
        \item Desglose secuencial de las operaciones térmicas y mecánicas (recepción, eviscerado, cocción, túneles de congelación) mediante diagramas de flujo de procesos (PFD).
        \item Catalogación técnica de la maquinaria industrial crítica (autoclaves, selladoras al vacío, sistemas de refrigeración por amoníaco).
        \item Mapeo teórico de los tiempos de ciclo y cuellos de botella característicos en la manipulación de materia prima perecedera.
    \end{itemize}

    \item \textbf{Fase 3: Mapeo de la Cadena de Suministro y Red Comercial.}
    \begin{itemize}[noitemsep, label=-]
        \item Rastreo documentado de los nodos de abastecimiento (\textit{upstream}): proveedores de diésel marino, mantenimiento de astilleros, hielo industrial y bobinas de hojalata/plástico.
        \item Identificación de los canales de distribución (\textit{downstream}): dinámica B2B con brokers de exportación y B2C con cadenas de supermercados, apoyándose en bases de datos aduaneras y registros comerciales.
    \end{itemize}

    \item \textbf{Fase 4: Parametrización de Variables y Consolidación Documental.}
    \begin{itemize}[noitemsep, label=-]
        \item Traducción de los procesos físicos investigados en elementos estadísticos, separando las variables del sistema en continuas (ej. temperatura de conservación, peso drenado) y discretas (ej. número de latas abolladas, defectos de sellado).
        \item Integración de los hallazgos en un documento marco que defina el universo poblacional, sirviendo como base teórica exacta para diseñar futuros modelos de muestreo.
    \end{itemize}


    
 % --------------------------------------------------------
% CRONOGRAMA NO LINEAL
% --------------------------------------------------------
\vspace{0.5cm}
\noindent \textbf{Cronograma de Trabajo Concurrente (No Lineal)} \\
El siguiente cronograma se diseña bajo un enfoque de ejecución concurrente. A diferencia de un modelo lineal puro en cascada, las actividades se superponen temporalmente para optimizar el tiempo de recopilación de información en campo y el cruce de datos.

\begin{table}[h!]
\centering
\renewcommand{\arraystretch}{1.5}
\begin{tabular}{|p{7.5cm}|c|c|c|c|c|}
\hline
\textbf{Actividad de Investigación} & \textbf{Sem 1} & \textbf{Sem 2} & \textbf{Sem 3} & \textbf{Sem 4} & \textbf{Sem 5} \\ \hline
Fase 1: Revisión Normativa y Estructural & \cellcolor{blue!20}X & \cellcolor{blue!20}X & & & \\ \hline
Fase 2: Levantamiento Analítico del Proceso & & \cellcolor{green!20}X & \cellcolor{green!20}X & \cellcolor{green!20}X & \\ \hline
Fase 3: Mapeo de la Cadena y Red Comercial & \cellcolor{orange!20}X & \cellcolor{orange!20}X & \cellcolor{orange!20}X & & \\ \hline
Fase 4: Parametrización y Documento Final & & & & \cellcolor{red!20}X & \cellcolor{red!20}X \\ \hline
\end{tabular}
\end{table}

% --------------------------------------------------------
% WEBGRAFÍA Y REFERENCIAS
% --------------------------------------------------------

\renewcommand{\refname}{Webgrafía y Referencias}
\begin{thebibliography}{9}

\bibitem{fao2022} 
Organización de las Naciones Unidas para la Alimentación y la Agricultura (FAO). (2022). 
\textit{El estado mundial de la pesca y la acuicultura: Hacia la transformación azul}. Recuperado de: \url{https://www.fao.org/fishery/es}

\bibitem{iso22000} 
Organización Internacional de Normalización (ISO). (2018). 
\textit{ISO 22000:2018 Sistemas de gestión de la inocuidad de los alimentos}. Recuperado de: \url{https://www.iso.org/standard/65464.html}


\bibitem{montgomery2019} 
Montgomery, D. C. (2019). 
\textit{Control estadístico de la calidad} (7ma ed.). Limusa Wiley. (Referencia base para la futura estructuración del muestreo de las variables).

\bibitem{bertalanffy} 
Von Bertalanffy, L. (1989). 
\textit{Teoría General de los Sistemas: Fundamentos, desarrollo, aplicaciones}. Fondo de Cultura Económica.

\end{thebibliography}   
    
    
    
\end{itemize}




\end{document}